%\documentclass[11pt,a4paper,twoside]{article}

\usepackage[utf8]{inputenc}

\usepackage[T1,T2A]{fontenc}

\usepackage[english.ancient,russian.ancient]{babel}

\usepackage{graphicx}

\usepackage{array}

\usepackage[doi]{samgtu-bib}

\usepackage{samgtu}

\usepackage{samgtu-name}

\usepackage{mathtools}

\mathtoolsset{showonlyrefs}

\usepackage{desclist}

\usepackage{euscript}

\usepackage{mathrsfs} 

\usepackage{multicol}

\usepackage{mathrsfs}

\usepackage{cite}

\usepackage{cmap}

\usepackage{qrcode}

\allowdisplaybreaks[4]

\usepackage[nodayofweek]{datetime}

\usepackage[thinlines]{easybmat}



\renewcommand{\JournalYear}{2018}

\renewcommand{\JournalVolume}{22}

\renewcommand{\JournalNumber}{4}



\newdate{JournalDateSign}{29}{12}{2018}% Дата подписания Вестника в печать

\renewcommand{\JournalDateSign}{\displaydate{JournalDateSign}}

\newdate{JournalDatePrint}{16}{01}{2019}% Дата выхода Вестника из печати

\renewcommand{\JournalDatePrint}{\displaydate{JournalDatePrint}}

%\renewcommand{\JournalUPL}{16.56} % Усл. печ. л.

%\renewcommand{\JournalUIL}{16.53} % Уч.-изд. л.

\renewcommand{\JournalUPL}{15.24} % Усл. печ. л.

\renewcommand{\JournalUIL}{15.21} % Уч.-изд. л.

\renewcommand{\JournalRegNumber}{220/18} % Регистрационный номер

\renewcommand{\JournalOrderNumber}{2} % Номер заказа



%\renewcommand{\JournalRMonth}{Март} % Месяц выхода Вестника

%\renewcommand{\JournalEMonth}{March} % Месяц выхода Вестника

%\renewcommand{\JournalRMonth}{Июнь} % Месяц выхода Вестника

%\renewcommand{\JournalEMonth}{June} % Месяц выхода Вестника

%\renewcommand{\JournalRMonth}{Сентябрь} % Месяц выхода Вестника

%\renewcommand{\JournalEMonth}{September} % Месяц выхода Вестника

\renewcommand{\JournalRMonth}{Декабрь} % Месяц выхода Вестника

\renewcommand{\JournalEMonth}{December} % Месяц выхода Вестника



\newcommand{\totop}[1]{\raisebox{6pt}[1pt][2pt]{#1}}

\newcommand{\tottop}[1]{\raisebox{12pt}[1pt][2pt]{#1}} 

\newcommand{\totttop}[1]{\raisebox{18pt}[1pt][2pt]{#1}} 

\newcommand{\tottttop}[1]{\raisebox{24pt}[1pt][2pt]{#1}} 

\newcommand{\totttttop}[1]{\raisebox{30pt}[1pt][2pt]{#1}} 

\newcommand{\tobot}[1]{\raisebox{-6pt}[1pt][2pt]{#1}} 

\newcommand{\tobbot}[1]{\raisebox{-12pt}[1pt][2pt]{#1}} 

\newcommand{\tobbbot}[1]{\raisebox{-18pt}[1pt][2pt]{#1}}



\graphicspath{{./00PIC/}}

%

\vsgtu{1569} %registration number

\FirstPage{236}

\LastPage{253}

\ResearchArticle

%

%\pdfcrop

%\draft

%

%\begin{document}

%

\hfill \qrcode[hyperlink,height=.8in]{http://doi.org/10.14498/vsgtu1569}

\vspace{-.8in}





\udc{517.982.45}

\msc{44A40, 35E20}



\rutitle{Построение операторного исчисления Микусинского~на~основе алгебры свертки обобщенных~функций. Решение задач математической~физики}

\rutitleshort{Построение операторного исчисления Микусинского\dots}



\entitle{Construction of Mikusinski operational calculus based on~the convolution algebra of distributions. Methods~for~solving mathematical physics problems}

\entitleshort{Construction of Mikusinski operational calculus\dots}



\ruauthor{И.~Л.~Коган}{К\,о\,г\,а\,н~И.\,Л.}

\enauthor{I.~L.~Kogan}{K\,o\,g\,a\,n~I.~L.}



\ruaffil{\ruaffid{5641}{Российский государственный аграрный университет -- МСХА им.~К.~А.~Тимирязева, 

\\ 

Россия, 127550, Москва, Тимирязевская ул., 49}.}



\enaffil{\enaffid{5641}{Russian State Agrarian University -- Moscow Agricultural Academy after K.~A.~Timiryazev, \\ 49,  Timiryazevskaya st., Moscow, 127550, Russian Federation}.}



\ruauthorsdetails{1}{% Количество авторов

\fullauthor{\rupid{34835}{Иосиф Леонидович Коган}}

\CAuthor % Автор-корреспондент

\orcid{0000-0002-1147-694X} % ORCID ID автора 

\regalia{кандидат технических наук} 

\position{доцент} 

\dept{каф. высшей математики} 

\email{ik_@list.ru}

}



\enauthorsdetails{1}{

\fullauthor{\enpid{34835}{Iosif L. Kogan}}

\CAuthor % Сorresponding Author

\orcid{0000-0002-1147-694X} % ORCID ID

\regalia{Cand. Techn. Sci.} 

\position{Associate Professor} 

\dept{Dept. of Higher Mathematics}

\email{ik_@list.ru} \vspace{-4mm}

}





\rukeywords{исчисление Микусинского, пространство обобщенных 

функций, свертка обобщенных функций, алгебра свертки, преобразование 

Лапласа, интеграл Дюамеля}

\enkeywords{calculus of Mikusi\'nski, space of distributions, 

convolution of distributions, convolution algebra, Laplace transform, 

Duhamel integral}



\ruabstract{Дается новое обоснование операторного исчисления Микусинского, целиком 

основанное на алгебре свертки обобщенных функций $D'_{+}$ и $D'_{-}$, 

применительно к решению линейных уравнений в частных производных с 

постоянными коэффициентами в области $(x;t)\in \mathbb R $ $( \mathbb R_{+} )\times \mathbb R_{+} $. 

Используемый математический аппарат основан на современном состоянии теории 

обобщенных функций, и~одним из основных его  отличий от теории Микусинского 

является то, что получаемые изображения являются аналитическими 

функциями комплексного переменного. Это позволяет в алгебре $D'_+ (x\in \mathbb R_{+})$ узаконить преобразование Лапласа, а~с~применением 

алгебры $D'_{-}$ распространить метод на область отрицательных значений 

аргумента. На классических примерах уравнений второго порядка 

гиперболического и параболического типа в случае $x\in \mathbb R$ излагаются вопросы 

определения фундаментальных решений и задачи Коши, а на отрезке и полупрямой 

$x\in \mathbb R_+ $~--- нестационарные задачи в~собственном смысле. Дается вывод 

общих формул для получения решения задачи Коши, а также схема определения 

фундаментальных решений операторным методом. При рассмотрении нестационарных 

задач приводится компактное доказательство теоремы Дюамеля и выведены 

формулы, позволяющие оптимизировать получение решений, в том числе 

с~разрывными начальными условиями. Для нахождения оригиналов приводятся 

примеры использования рядов сверточных операторов обобщенных функций. 

Предложенный подход по сравнению с классическим операционным исчислением, 

основанным на преобразовании Лапласа, и~теорией Микусинского, обладая для 

обычных функций одинаковыми соотношениями <<оригинал-изображение>> на 

положительной полуоси, позволяет рассматривать уравнения, заданные на всей 

оси, упростить получение и форму представления решений. Приведенные примеры 

иллюстрируют возможности и дают оценку эффективности использования 

операторного исчисления.}



\enabstract{A new justification is given for the Mikusinsky operator calculus entirely based on the convolution algebra of generalized functions $D'_+$ and $D'_-$, as applied to the solution of linear partial differential equations with constant coefficients in the region $(x;t)\in \mathbb R (\mathbb R_{+})\times \mathbb R_{+}$. The mathematical apparatus used is based on the current state of the theory of generalized functions and its one of the main differences from the theory of Mikusinsky is that the resulting images are analytical functions of a complex variable. This allows us to legitimate the Laplace transform in the algebra $D'_{+} $ $( x\in \mathbb R_{+} )$, and apply the algebra to the region of negative values of the argument with the use of algebra $D'_{-}$. On classical examples of second-order equations of hyperbolic and parabolic type, in the case $x\in \mathbb R$, questions of the definition of fundamental solutions and the Cauchy problem are stated, and on the segment and the half-line $x\in \mathbb R_{+}$, non-stationary problems in the proper sense are considered. We derive general formulas for the Cauchy problem, as well as circuit of fundamental solutions definition by operator method. When considering non-stationary problems we introduce the compact proof of Duhamel theorem and derive the formulas which allow optimizing obtaining of solutions, including problems with discontinuous initial conditions. Examples of using series of convolution operators of generalized functions are given to find the originals. The proposed approach is compared with classical operational calculus based on the Laplace transform, and the theory of Mikusinsky, having the same ratios of the original image on the positive half-axis for normal functions allows us to consider the equations posed on the whole axis, to facilitate the obtaining and presentation of solutions. These examples illustrate the possibilities and give an assessment of the efficiency of the use of operator calculus.}



\newdate{kogana}{17}{10}{2017}

\date{\displaydate{kogana}}

\newdate{koganb}{11}{02}{2018}

\revisiondate{\displaydate{koganb}}

\newdate{koganc}{12}{03}{2018}

\accepteddate{\displaydate{koganc}}



\newdate{kogane}{28}{03}{2018}

\renewcommand{\JournalDataOnline}{\displaydate{kogane}}



%\ForPeerReview

\makerutitle





\Section[n]{Введение} Операторное, или операционное, исчисление достаточно давно 

и успешно применяется в математической физике, главным образом при решении 

нестационарных задач. Исторически первой практической реализацией этого 

исчисления является символический метод Хевисайда, который ввел правила 

обращения с операторами дифференцирования и интегрирования, рассматриваемыми 

как алгебраические величины~\cite{kogan:bib1}. К удобству такого подхода следует отнести 

то, что формальные вычисления можно отделить от математического содержания 

задачи. 



Строгое обоснование метода Хевисайда было дано с помощью 

интегрального преобразования Лапласа на основании теории функций 

комплексного переменного~[\citen{kogan:bib1}--\citen{kogan:bib3}]. 

Подход к операционному исчислению как 

приложению этой теории "--- классическое операционное исчисление доминирует в 

современной учебной и научно"=технической литературе, например~[\citen{kogan:bib4}--\citen{kogan:bib6}]. 



Введение обобщенных функций существенно расширяет возможности операционного 

исчисления. Наиболее распространенной является предложенная  Л.~Шварцем и Ж.-Л.~Лионсом 

теория преобразования Лапласа для обобщенных функций медленного роста, 

заданных на положительной полуоси и~образующих сверточную алгебру $S'_{+} 

\subset D'_{+}$~[\citen{kogan:bib7}--\citen{kogan:bib13}]. К удобствам теории можно отнести то, что оригинал 

в смысле классического операционного исчисления является и оригиналом в 

обобщенном смысле, при этом сразу же получаются изображения дельта"=функции 

Дирака и ее производных. Кроме того, связь между оригиналом и изображением, 

а также основные теоремы совпадают с~классическим операционным исчислением. 



Одним из ограничений применения классического операционного исчисления и 

преобразования Лапласа обобщенных функций медленного роста, накладываемых на 

используемые функции с не ограниченным справа носителем, является то, что их 

рост в бесконечности не должен превышать экспоненциальный. 

Возможность  в  некоторых случаях обойти указанное ограничение позволяет операторный метод 

Микусинского~\cite{kogan:bib14}, представляющий возврат на более высоком уровне к 

первоначальной операторной точке зрения, основываясь на введении

коммутативного кольца операторов, под которыми понимаются числовые 

операторы, функции, заданные на положительной полуоси, и собственно 

операторы. 

В этом кольце определены линейные операции сложения и умножения 

на числа, а в качестве операции умножения используется свертка функций 

(операторов), единичным элементом для которой является дельта"=функция 

Дирака. В своей теории Микусинский исключил преобразование Лапласа и тем 

самым снял ограничение на рост функций в бесконечности. Получаемые здесь 

изображения являются функциями от оператора дифференцирования, 

рассматриваемого, аналогично исчислению Хевисайда, как алгебраическая 

величина. 

Между операторным исчислением Микусинского и классическим 

операционным исчислением существует формальное сходство: изображения функций 

по Лапласу совпадают с соответствующими изображениями по Микусинскому и 

между этими методами можно установить изоморфизм. 



Дальнейшее развитие идей 

Микусинского дано в работах~\cite{kogan:bib15,kogan:bib16}, в которых операторное исчисление 

строится на более широкой основе "--- алгебре свертки обобщенных функций 

$D'_{+}$ и $D'_{-}$, а связь между оригиналом и изображением находится из 

заданного сверточного соотношения. Это дает возможность не рассматривать 

поведение функций в бесконечности и распространить метод на область 

отрицательных значений аргумента. 

Другим отличием от теории Микусинского 

является то, что получаемые изображения являются аналитическими функциями 

комплексного переменного; при этом, как и в преобразовании Лапласа 

обобщенных функций медленного роста, комплексная переменная есть изображение 

сверточного оператора дифференцирования. Такое построение теории, сохраняя 

преимущества, заложенные в~\cite{kogan:bib14}, позволяет находить решения на всей оси, 

применять методы теории функций комплексного переменного и рассматривать 

классическое операционное исчисление как свою составную часть. Кроме того, 

использование алгебры свертки обобщенных функций расширяет возможности 

операторного исчисления.



Настоящая работа посвящена приложению операторного исчисления~\cite{kogan:bib15,kogan:bib16} к 

решению некоторых задач математической физики одного пространственного переменного $x$ и временной координаты $t$. 

Рассматриваются уравнения с постоянными коэффициентами

\begin{equation}

\label{kogan:eq1}

L\Bigl( \frac{\partial }{\partial x}; \frac{\partial}{\partial t} \Bigr) u(x;t)\equiv a_1 u_{tt} +b_1 u_t +a_2 u_{xx} +b_2 u_x + c u=f(x;t),

\end{equation} 

где $x\in \mathbb R$ или $ \mathbb R_+$, $t\in \mathbb R_+$. 

Основные случаи уравнения \eqref{kogan:eq1} при $a_2<0$:

\begin{equation}

\label{kogan:eq2}

\begin{aligned}

& a_1 =1 \quad \text{(гиперболический случай)},\\

& a_1 =0;\quad b_1 =1 \quad \text{(параболический случай)}.

\end{aligned}

\end{equation}



Выделяются две задачи.



\begin{enumerate}

\item[1.] Задача Коши, когда $x\in \mathbb R$, с начальными условиями

\begin{equation}

\label{kogan:eq3}

\begin{aligned}

& u(x;0)=u_0 (x), \\

& u_t (x;0)=u_1 (x) \quad \text{(гиперболический случай)}. 

\end{aligned} 

\end{equation}



\item[2.] Нестационарная задача в собственном смысле, когда $x \! \in \! \mathbb R_+$ или 

${x \! \in \! [0; l]}$ и рассматривается однородное уравнение \eqref{kogan:eq1}, \eqref{kogan:eq2}, 

т.~е. $f(x;t)=0$, при нулевых начальных условиях Коши \eqref{kogan:eq3}. 

В этом заключается  собственный смысл задачи. 

В случае $x\in \mathbb R_+$ будут использованы следующие  граничные условия: 

\begin{gather}

\label{kogan:eq4}

u(0;t)=g(t),

\quad

\lim \limits_{x\to \infty } u(x;t)=0;\\

\label{kogan:eq5}

u_x (0;t)=v(t),

\quad

\lim \limits_{x\to \infty } u(x;t)=0.

\end{gather}

Если $x\in [0;l]$, то наряду с указанными условиями, когда $x=0$, задаются условия при $x=l$.



\end{enumerate}



Указанный выбор обусловлен тем, что первый класс задач представляет реализацию 

решения на всей полуплоскости $(x;t)\in \mathbb R\times \mathbb R_+$, что невозможно в~рамках классического операционного исчисления и теории Микусинского, а при 

решении второго класса задач, типичных для применения операционного 

исчисления, имеется возможность оптимизировать как само получение решения, 

так и форму его представления. Приводимая методика и теоремы могут быть 

обобщены для решений уравнений более высокого порядка.



\smallskip



\Section{Задача Коши $\boldsymbol{(x\in \mathbb R)}$} Решение ищется в алгебре 

свертки $D'_{+} (t)\hm \times \bigl(D'_+ (x)\cup D'_- (x)\bigr)$ с единичным элементом 

$\delta (x;t)=\delta (x)\times \delta (t)$ "--- дельта"=функцией Дирака. 

Операторами дифференцирования $n$"=ного порядка по переменным $x$ и $t$ в~смысле теории 

обобщенных функций здесь будут $\bigl(\delta'(x)\bigr)^n$ и $\bigl(\delta'(t)\bigr)^n$, 

причем $(\delta')^0=\delta$, и уравнение \eqref{kogan:eq1} примет вид

\begin{equation}

\label{kogan:eq6}

L\big(\delta'(x);\delta'(t) \big)\ast u(x;t)=\vartheta(t) f(x;t).

\end{equation}



Следуя методологии Л.~Шварца для обыкновенных дифференциальных уравнений с 

постоянными коэффициентами~\cite{kogan:bib9}, а также используя известную теорему о существовании фундаментального решения 

у любого линейного дифференциального оператора с постоянными коэффициентами 

в $D'$, например~\cite{kogan:bib17}, можно сформулировать следующее утверждение.



\smallskip



\hypertarget{kogan:th1}{}



\begin{theorem}[1]Оператор в левой части \eqref{kogan:eq6} обратим, т.~е. фундаментальное решение уравнения \eqref{kogan:eq1} или \eqref{kogan:eq6} в смысле алгебры свертки имеет вид

\begin{equation}

\label{kogan:eq7}

\left(L\bigl(\delta'(x);\delta'(t) \bigr) \right)^{-1}=\vartheta (t)E_0 (x;t),

\end{equation}

где $E_0(x;t)$ "--- решение соответствующего однородного уравнения, отвечающее начальным условиям

\begin{equation}

\label{kogan:eq8}

\begin{aligned}

& E_0 (x;0)=\delta (x) \quad \text{$($параболический случай$)$}, \\

& E_0 (x;0)=0; \quad (E_0)_t (x;0) = \delta (x) \quad \text{$($гиперболический случай$)$}. 

\end{aligned}

\end{equation}

\end{theorem}



\smallskip



\begin{proof} Доказательство приведем для гиперболического \linebreak случая. 

Применяя операторы дифференцирования 

по $t$ с учетом начальных условий \eqref{kogan:eq8}, имеем

\begin{equation}

\label{kogan:eq9} 

\begin{aligned}

&\delta'(t) \ast \vartheta (t)E_0 (x;t)=\vartheta (t)(E_0)_t (x;t)+\delta 

(t)E_0 (x;0)=\vartheta (t)(E_0)_t (x;t);\\

&\big(\delta'(t) \big)^2 \ast \vartheta (t) E_0 (x;t)  =\vartheta (t)(E_0)_{tt} 

(x;t)+\delta (t)(E_0)_t (x;0) = \\

& \hspace{7cm} =\vartheta (t)(E_0)_{tt} (x;t)+\delta (x;t).

\end{aligned}

\end{equation}

Отсюда, учитывая, что $E_0 (x;t)$ "--- решение однородного уравнения, получаем 

\begin{equation}

L\big(\delta'(x);\delta'(t) \big)\ast \vartheta (t)E_0 (x;t)=\vartheta 

(t)L\big(\delta'(x);\delta'(t) \big) \ast E_0 (x;t)+\delta (x;t)=\delta (x;t).

\end{equation} 







\hfill\end{proof}



\smallskip



\begin{remark}[1] Теорема~\hyperlink{kogan:th1}{1}, как нетрудно видеть, справедлива для оператора $L$ 

типа \eqref{kogan:eq6} любого порядка относительно $t$, если коэффициент при $\big(\delta 

'(t)\big)^n$ максимального порядка равен единице, а начальные условия Коши 

\begin{equation}

\label{kogan:eq10}

\begin{aligned}

&E_0 (x;0)=(E_0)_t (x;0)=\ldots=\frac{\partial^{(n-2)}E_0 (x;0)}{ \partial t^{(n-2)}}=0;\\

& \frac{\partial^{(n-1)}E_0 (x;0) }{ \partial t^{(n-1)} }=\delta (x).

\end{aligned}

\end{equation}

Функция $E_0 (x;t)$ называется {\it фундаментальным решением задачи Коши однородного уравнения}~\cite{kogan:bib18}. 

Таким 

образом, теорема~\hyperlink{kogan:th1}{1} определяет связь между фундаментальным решением 

неоднородного уравнения в алгебре свертки и фундаментальным решением задачи 

Коши соответствующего однородного уравнения. 

\end{remark}



\smallskip



Дальнейшее изложение  основывается на следующей лемме.



\smallskip



\begin{lemma}Если $k\le n-1,$ где $n$ "--- порядок оператора $L$ относительно $t,$ и выполняется хотя бы одно из двух условий{\/\rm:} 

\begin{enumerate}

\item[\rm 1)] фундаментальное решение $E(x;t)=\vartheta (t)E_0 (x;t)$ при любом фиксированном $t$ финитно{\rm;}

\item[\rm 2)] $e_k (x)$ "--- финитная функция$,$

\end{enumerate}

то справедливо сверточное равенство

\begin{equation}

\label{kogan:eq11}

E(x;t)\ast e_k (x)\big(\delta'(t)\big)^k=\vartheta (t)\frac{\partial^kE_0 (x;t)}{\partial t^k}\ast e_k (x),

\end{equation}

где дифференцирование понимается в смысле теории обобщенных функций\rm.

\end{lemma}





\smallskip



\begin{proof}

Рассмотрим случай $k=0$. Оставим общий индекс $k$ у функции $e_k(x)$. 

Нужно показать, что

\begin{equation}

\label{kogan:eq12}

E(x;t) \ast e_k (x)\delta (t)=E(x;t) \ast e_k (x).

\end{equation}



Действительно, по определению свертки и прямого произведения обобщённых 

функций для $\forall \varphi (x;t)\in D(x;t)$ имеем

\begin{multline*}

E(x;t) \ast e_k (x) \delta (t)=\big(E(x;t)\times e_k (y)\times \delta (z),\varphi (x+y;t+z) \big)=\\

=\left(E(x;t),\big(e_k (y),(\delta (z),\varphi (x+y;t+z))\big) \right)=\\

=\big(E(x;t), \big(e_k (y),\varphi  (x+y;t)\big) \big)=E(x;t)\ast e_k (x),

\end{multline*}

так как $\varphi (x;t)\in D(x)$ при любом $t=t_0$, а свертка $E(x;t_0 )\ast 

e_k (x)$, согласно условию леммы, всегда определена. 



Теперь перейдем к  общему случаю. Исходя из правила дифференцирования свертки функций, учитывая 

условие~\eqref{kogan:eq10} и соотношение \eqref{kogan:eq12}, имеем

\[

E(x;t)\ast e_k (x)\big(\delta'(t) \big)^k=\frac{\partial ^k\big(\vartheta (t)E_0 (x;t) \big)}{\partial t^k}\ast e_k (x)\delta (t)=\vartheta (t)\frac{\partial 

^k E_0 (x;t)}{\partial t^k}\ast e_k (x).

\]



\hfill

\end{proof}



\smallskip



\begin{theorem}[2]Решение задачи Коши \eqref{kogan:eq1}--\eqref{kogan:eq3} в алгебре свертки имеет вид

\begin{equation}

\label{kogan:eq13} 

\vartheta (t)u(x;t)=\vartheta (t)E_0 (x;t)\ast \vartheta (t)f(x;t)+\vartheta 

(t)\sum\limits_{k=0}^{n-1} \frac{\partial^k E_0 (x;t)}{\partial t^k} \ast e_k (x),

\end{equation}

где $\vartheta (t)E_0 (x;t)$ "--- фундаментальное решение, $n$ "--- порядок оператора $L$ по переменной $t$, $e_k (x)$ определяются из выражения $L$ и начальных условий. 



В~рассматриваемых случаях:

\begin{equation}

\label{kogan:eq14}

\begin{aligned}

& n=1 \text{ $($параболический случай$){:}$ } e_0(x) = u_0(x);\\

& n=2 \text{ $($гиперболический случай$){:}$ } e_0 (x)=b_1 u_0 (x)+u_1 (x);\; e_1(x) =u_0 (x). 

\end{aligned}\!\!

\end{equation}

\end{theorem}



\smallskip



\begin{proof}

Пусть $u(x;t)$ "--- решение задачи \eqref{kogan:eq1}--\eqref{kogan:eq3}. 

Применяя процедуру \eqref{kogan:eq9} к~$\vartheta (t)u(x;t)$ с учетом начальных условий \eqref{kogan:eq3}, получаем

\begin{equation}

\label{kogan:eq15}

L\big(\delta'(x);\delta'(t) \big) \ast \vartheta (t)u(x;t)=\vartheta 

(t)f(x;t)+\sum\limits_{k=0}^{n-1} e_k (x)\big(\delta'(t) \big)^k,

\end{equation}

где $e_k (x)$ находятся согласно \eqref{kogan:eq14}. После применения оператора \eqref{kogan:eq7} к обеим частям 

уравнения \eqref{kogan:eq15} с учетом леммы следует утверждение теоремы \eqref{kogan:eq13},~\eqref{kogan:eq14}. 



Отметим, что начальные условия \eqref{kogan:eq3} входят в правую часть \eqref{kogan:eq14}, а само 

решение этого уравнения по построению отвечает заданным начальным условиям 

Коши. Так как $t\ge 0$,  множитель $\vartheta (t)$ в формуле \eqref{kogan:eq13} можно 

опустить.

\end{proof}



\smallskip



{\bf 1.1. Схема определения фундаментального решение операторным методом.}

Применим последовательно операторные преобразования по $t$ и $x$ к выражению 

фундаментального решения \eqref{kogan:eq7}. Преобразование по $t$ в алгебре $D'_+$ с 

единичным элементом $\delta (t)$ (при этом переменная $x$ рассматривается как параметр) 

дает 

\begin{equation}

\label{kogan:eq16}

E(x;t)=\vartheta (t)E_0 (x;t)\leftrightarrow E(x;p);

\quad

\big(\delta'(t) \big)^k \leftrightarrow p^k \quad (k=0,1,2,\ldots).

\end{equation}



Преобразование по $x$ в алгебрах $D'_+$ и $D'_-$ с единичным элементом 

$\delta (x)$, где комплексная переменная $p$ рассматривается как параметр, дает 

\begin{equation}

\label{kogan:eq17}

E(x;p)\leftrightarrow E(p_1 ;p);

\quad

\big(\delta'(x) \big)^k \leftrightarrow p_1^k

\quad

(k=0,1,2,\ldots).

\end{equation}

Следовательно, изображение фундаментального 

решения \eqref{kogan:eq7} согласно \eqref{kogan:eq16} и~\eqref{kogan:eq17} будет иметь вид

\begin{equation}

\label{kogan:eq18}

E\big(p_1 ;p\big)=\bigl(L(p_1 ;p) \bigr)^{-1}.

\end{equation}



Отсюда следует правило: {\sl чтобы найти изображение фундаментального решения \eqref{kogan:eq18}, нужно в многочлене $L$ оператора \eqref{kogan:eq1} символы $\partial 

/\partial x$ и $\partial/\partial t$ заменить соответственно на $p_1$ и $p$, а потом взять обратную величину от полученного выражения.} 

Оригинал $E(x;t)$ находится обратными преобразованиями по схеме

\begin{equation}

\label{kogan:eq19}

E (p_1 ;p ) \to \begin{matrix}

E_+ (p_1 ;p) \\

E_- (p_1 ;p) 

\end{matrix} \xrightarrow{D'_+;D'_-}  

\begin{matrix}

E(x_+;p) \\

E(x_-;p)

\end{matrix}  \xrightarrow{D'_+} E(x;t).

\end{equation}



Чтобы выделить оригиналы для $x\in \mathbb R_+$ и $x\in \mathbb  R_-$, требуется найти 

составляющие $E(p_1;p )$: $E_+ (p_1,p ) \in D'_+$ и $E_- (p_1 ,p ) \in D'_-$.

Для этого нужно провести разложение этого изображения на элементарные 

дроби по переменной $p_1$ и~воспользоваться предельными соотношениями при 

$\mathop{\rm Re} p_1 \to \pm \infty$, учитывая при этом, что одновременно с $\mathop{\rm Re} p\to 

+\infty$ изображение должно стремиться к нулю. Напомним, что в алгебрах 

$D'_+$ и $D'_-$ изображения находятся в разных полуплоскостях от мнимой 

оси, а оригиналы для одинаковых изображений в~этих алгебрах отличаются лишь 

знаком~\cite{kogan:bib15,kogan:bib16}. 



Впрочем, можно, ограничившись алгеброй $D'_+$, найти оригинал для случая 

$x \in \mathbb R_+$, а для случая $x\in \mathbb R_-$ оригинал определить из соображений 

симметрии.



\smallskip



{\bf 1.2. Решения задачи Коши операторным методом в области \linebreak $\boldsymbol{\mathbb R_+ (t) \times} \boldsymbol{\mathbb  R(x)}$ на классических примерах уравнения теплопроводности и 

волнового уравнения.} Для уравнения теплопроводности 

(параболического типа) задача Коши имеет вид 

\begin{equation}

\label{kogan:eq20}

\Bigl( \frac{\partial }{\partial t}-a^2\frac{\partial ^2}{\partial x^2} \Bigr)u(x;t)=f(x;t),

\quad

a>0;

\quad

u(x;0)=u_0 (x).

\end{equation}

Решение этой задачи на основании \eqref{kogan:eq13}, \eqref{kogan:eq14} ищем в 

виде 

\begin{equation}

\label{kogan:eq21}

u(x;t)=\vartheta (t)E_0 (x;t) \ast \vartheta (t)f(x;t)+\vartheta (t)E_0 

(x;t) \ast u_0 (x).

\end{equation}

Определяем фундаментальное решение $E(x;t)=\vartheta 

(t)E_0 (x;t)$. Согласно правилу \eqref{kogan:eq18} и схеме \eqref{kogan:eq19}:

\[

E(p_1 ;p)=(p-a^2p_1^2 )^{-1}=E_+ (p_1 ;p )+E_- (p_1 ;p),

\]

где 

$$

E_+ (p_1  ;p)=\frac{1}{2a\sqrt p }\Bigl( \frac{1}{p_1 +{\sqrt p } / a} \Bigr),

\quad 

E_- (p_1 ;p)=-\frac{1}{2a\sqrt p} \Bigl( \frac{1}{p_1 -{\sqrt p} / a}\Bigr).

$$

Используя известное соотношение 

преобразования Лапласа~\cite{kogan:bib19}

$$

\frac{1}{\sqrt p }e^{-k\sqrt p } \leftrightarrow \frac{1}{\sqrt {\pi t}}e^{-\frac{k^2}{4t}} \quad  

(k>0),

$$ 

последовательно получаем: 

$$

E_+ (p_1 ;p)\leftrightarrow  \frac{1}{2a\sqrt p }e^{-\frac{x}{a}\sqrt p }\leftrightarrow \frac{\vartheta (t)}{2a\sqrt {\pi t}  }e^{-\frac{x^2}{4a^2t}}, \quad x\in \mathbb R_+ ;

$$

$$

E_- (p_1 ;p)\leftrightarrow \frac{1}{2a\sqrt p }e^{-\frac{|x|}{a}\sqrt p }\leftrightarrow \frac{\vartheta (t)}{2a\sqrt {\pi t}  }e^{-\frac{x^2}{4a^2t}},\quad x\in \mathbb R_-.

$$ 

Таким образом, 

\begin{equation}

\label{kogan:eq22}

E(x;t)=\frac{\vartheta (t)}{2a\sqrt {\pi t}}e^{-\frac{x^2}{4a^2t}},

\end{equation}

где 

$x\in \mathbb R$. 



На основании \eqref{kogan:eq21}, \eqref{kogan:eq22} получаем классическую формулу Пуассона:

\begin{equation}

\label{kogan:eq23}

u(x;t)\! = \!\int_0^t \int_{-\infty }^{+\infty } \frac{f(\xi;\tau)}{2a\sqrt{ 

\pi (t-\tau )} }  e^{-\frac{(x-\xi )^2}{4a^2(t-\tau )}} d\xi d\tau 

\!+\! \frac{\vartheta (t)}{2a\sqrt {\pi t} }\int_{-\infty}^{+\infty } u_0 (\xi)e^{-\frac{(x-\xi)^2}{4a^2t}}d\xi.

\end{equation}





Для волнового уравнения (гиперболического вида) задача Коши имеет вид

\begin{equation}

\label{kogan:eq24}

\begin{aligned}

& \Bigl( \frac{\partial ^2}{\partial t^2}-a^2\frac{\partial ^2}{\partial 

x^2} \Bigr)u(x;t)=f(x;t),\\

& a>0;

\quad

u(x;0)=u_0 (x),

\quad

u_t (x;0)=u_1 (x).

\end{aligned}

\end{equation}

 

 

Решение согласно \eqref{kogan:eq13}, \eqref{kogan:eq14} ищем в виде

\begin{multline}

\label{kogan:eq25}

u(x;t)=\vartheta (t)E_0 (x;t)\ast \vartheta (t)f(x;t) +

\\ +\vartheta (t)E_0 

(x;t)\ast u_1 (x)

+\vartheta (t)\frac{\partial E_0 (x;t)}{\partial t}\ast u_0 

(x).

\end{multline}



Определим фундаментальное решение. 

На основании \eqref{kogan:eq18}, \eqref{kogan:eq19}, получаем 

$$

E(p_1 ;p)=(p^2-a^2p_1^2 )^{-1}=E_+ (p_1 ;p)+E_- (p_1 ;p),

$$

где 

$$

E_+  (p_1;p )=\frac{1}{2ap(p_1 +p / a )}, \quad 

E_- (p_1 ;p )=-\frac{1}{2ap(p_1 -p/a )}.

$$ 

Используя теоремы подобия и~запаздывания оригинала, имеем

$$

E_+ (p_1 ;p ) \leftrightarrow \frac{1}{2a^2(\frac{p}{a} )}e^{-x (\frac{p}{a} )} \leftrightarrow 

\frac{1}{2a}\vartheta (at-x), \quad x\in \mathbb R_+;

$$

$$

E_- (p_1 ;p )  \leftrightarrow \frac{1}{2ap}e^{x(\frac{p}{a} )} \leftrightarrow 

\frac{1}{2a}\vartheta (at+x),\quad  x\in \mathbb R_-.

$$

Таким образом, 

\begin{equation}

\label{kogan:eq26}

E(x;t)=\frac{1}{2a}\vartheta (at- |  x |  )=\begin{cases}

1/2a, & \text{если} \quad  |  x |  <at,\\

~~~0, &\text{если} \quad  |  x |  >at, 

\end{cases}

\end{equation}

где $x\in  \mathbb R.$



Принимая во внимание четность дельта-функции Дирака и равенство, следующее из \eqref{kogan:eq26},

$$\frac{\partial E(x;t)}{\partial t}=\frac{1}{2}\delta 

(at- |  x |  ),$$ получаем 

$$

\frac{\partial 

E(x;t)}{\partial t}\ast u_0 (x)=\frac{1}{2}\bigl( u_0 (x+at)+u_0 (x-at) \bigr).

$$ 

Отсюда и на основании \eqref{kogan:eq25}, \eqref{kogan:eq26} после элементарных преобразований

получаем классическую формулу Даламбера:

\begin{multline}

\label{kogan:eq27}

u(x;t)=\frac{1}{2a}\int_0^t \int_{x-a(t-\tau )}^{x+a(t-\tau )} f(\xi ;\tau) d\xi d\tau +\\

+\frac{1}{2a}\int _{x-at}^{x+at} u_1 (\xi )d\xi 

+\frac{1}{2}\bigl(u_0 (x+at)+u_0 (x-at)\bigr).

\end{multline}

%Учитывая, что $\frac{\partial E(x;t)}{\partial t}=\frac{1}{2}\delta 

%(at- |  x |  )$ и четность дельта функции Дирака, $\frac{\partial 

%E(x;t)}{\partial t}\ast u_0 (x)=\frac{1}{2}\big[ u_0 (x+at)+u_0 (x-at)\big]$. Отсюда и на основании 

%\eqref{kogan:eq25}, \eqref{kogan:eq26} после элементарных преобразований получаем классическую формулу Даламбера

%\begin{multline}

%\label{kogan:eq27}

%u(x;t)=\frac{1}{2a}\int_0^t \int_{x-a(t-\tau )}^{x+a(t-\tau )} f(\xi;\tau) d\xi d\tau +\\

%+\frac{1}{2a}\int \limits_{x-at}^{x+at} u_1 (\xi ) d\xi 

%+\frac{1}{2}\big[ u_0 (x+at)+u_0 (x-at) \big].

%\end{multline}



\smallskip



\Section{Нестационарная задача в собственном смысле, $\boldsymbol{ x \!\in\! \mathbb R_+}$ $\boldsymbol{(x\!\in\! [0;l] )}$} Здесь используется 

иной подход, чем в первом типе задач. Применим операторное преобразование  по $t$ в алгебре $D'_+$ к исходному однородному 

уравнению \eqref{kogan:eq1}, \eqref{kogan:eq2} при нулевых начальных условиях Коши \eqref{kogan:eq3} и заданных 

граничных условиях. 

С учетом теоремы о дифференцировании по параметру  соотношения <<оригинал"=изображение>>~\cite{kogan:bib16} имеем

\begin{equation}

\label{kogan:eq28}

u(x;t)\leftrightarrow U(x;p);

\quad

\frac{\partial^ ku(x;t)}{\partial t^k}\leftrightarrow p^kU(x;p);

\quad

\frac{\partial ^ku(x;t)}{\partial x^k}\leftrightarrow 

\frac{d^kU(x;p)}{dx^k}.

\end{equation}



Далее решается операторное уравнение относительно $U(x;p)$, т.~е. краевая 

задача для обыкновенного дифференциального уравнения, и по найденному 

изображению находится оригинал $u(x;t)$. 

Наряду со стандартными приемами 

классического операционного исчисления для определения $u(x;t)$ бывает более 

удобным использовать методы алгебры свертки $D'_+$, разлагая $U(x;p)$ в ряд,

и рассматривать ряды сверточных операторов обобщенных функций. 

В~качестве 

примера приведем доказательство двух соотношений, представленных в виде, 

удобном для использования в~дальнейшем. 



\smallskip



\begin{newthm}{Предложение} Пусть $f(t)\leftrightarrow F(p)$ и $c,$ $d\in \mathbb R_+$ "--- постоянные$,$ тогда

\begin{equation}

\label{kogan:eq29}

\begin{aligned}

&\frac{F(p)e^{-cxp}}{1+e^{-dp}}\leftrightarrow 

f(t-cx)+\sum\limits_{k=1}^\infty (-1)^kf(t-cx-kd) =

\\ 

& \hspace{6.5cm} =\sum\limits_{k=0}^\infty (-1)^kf( t-cx-kd);\\

&\frac{F(p)e^{-cxp}}{1-e^{-dp}}\leftrightarrow 

f(t-cx)+\sum\limits_{k=1}^\infty f(t-cx-kd)=\sum\limits_{k=0}^\infty f(t-cx-kd).

\end{aligned}

\end{equation}

\end{newthm}



\smallskip



\begin{proof} Ограничимся доказательством первого соотношения. Используя разложение в 

ряд и теоремы умножения изображений и запаздывания оригинала, приходим к 

выражению, представляющему собой свертку функции с бесконечным рядом операторов сдвига:

\begin{multline}

\label{kogan:eq30}

\frac{F(p)e^{-cxp}}{1+e^{-dp}}=F(p)e^{-cxp}

\biggl(1+\sum\limits_{k=1}^\infty (-1)^k e^{-kdp} \biggr)\leftrightarrow \\

\leftrightarrow f(t-cx)\ast 

\biggl( \delta (t)+\sum\limits_{k=1}^\infty (-1)^k\delta  (t-kd) \biggr)\leftrightarrow \\

\leftrightarrow f(t-cx) + \sum\limits_{k=1}^\infty (-1)^kf(t-cx-kd) 

\end{multline}

или 

\begin{equation}

\frac{F(p)e^{-cxp}}{1+e^{-dp}}\leftrightarrow \sum\limits_{k=0}^\infty 

(-1)^kf(t-cx-kd).

\end{equation}



\hfill

\end{proof}



Учитывая неравенство $t\ge cx+kd$, отметим, что при любом значении $t$ ряды 

\eqref{kogan:eq29} содержат конечное число слагаемых. 



Далее на примере волнового уравнения 

рассмотрим решения задач, когда $x\in \mathbb R_+$ и $x\in [0;l]$. 



\smallskip



{\bf 2.1. Примеры граничных задач для волнового уравнения.} Рассмотрим 

сначала случай $x\in \mathbb  R_+$. 

Пусть начало струны $x=0$ свободно и~по 

неподвижной струне наносится воздействие такое, что выполняется условие 

$u_x  (0;t)=v(t)$. 

Таким образом, решаем волновое уравнение при нулевых начальных 

условиях Коши \eqref{kogan:eq3} и граничных условиях \eqref{kogan:eq5}: 

\begin{equation}

\label{kogan:eq30}

\begin{aligned}

& \Bigl( \frac{\partial ^2}{\partial t^2}-a^2\frac{\partial ^2}{\partial 

x^2} \Bigr)u(x;t)=0;

\quad

u(x;0)=u_t (x;0)=0;\\

& u_x (0;t)=v(t);

\quad

\lim _{x\to \infty } u(x;t)=0.

\end{aligned}

\end{equation}

Изображение 

$$U(x;p)=-\frac{aV(p)}{p}e^{-\frac{x}{a}p}$$

"--- решение операторного уравнения 

$$

p^2U-a^2\frac{d^2U}{dx^2}=0,$$  соответствующего условиям  \eqref{kogan:eq30},  отвечающее

граничным условиям 

$$\frac{dU}{dx}(0;p)=V(p), \quad \lim \limits_{x\to \infty, \mathop{\rm Re}P>0} U(x;p)\to 0.

$$ Отсюда 

\begin{equation}

\label{kogan:eq31}

u(x;t)=-av(t)\ast \vartheta \Bigl( t-\frac{x}{a}\Bigr)

\end{equation}

"--- решение задачи в алгебре свертки.



Выпишем решения для частных случаев функции $v(t)$ в начальных условиях:



\begin{itemize}

\item[а)] если  $v(t)$ "--- локально интегрируемая функция, то 

$$u(x;t)=-aV\Bigl( t-\frac{x}{a} \Bigr),$$ где 

$$V(z)=\int_0^z v(\tau )d\tau;$$ 



\item[б)] если  $v(t)=\delta (t)$, т.е. по струне наносится мгновенный удар, то

$$

u(x;t)=-a\vartheta \Bigl(t-\frac{x}{a}\Bigr).

$$





	

\end{itemize}



Пусть теперь $x\in [0;l]$. Будем решать такую же задачу для 

уравнения \eqref{kogan:eq30}, но для закрепленной на конце $x=l$ струны, 

т.~е. при нулевых начальных условиях Коши будут выполняться граничные условия:

\begin{equation}

\label{kogan:eq311}

u_x (0;t)=v(t),

\quad

u(l;t)=0.

\end{equation}



В этом случае, учитывая разложение 

\eqref{kogan:eq29}, получаем

$$

U(x;p)=-a\frac{V(p)}{p}\left(  \frac{e^{-\frac{x}{a}p}-e^{-(\frac{2l-x}{a})p}}{1+e^{-\frac{2l}{a}p}} \right)

$$

"--- решение операторного уравнения,  отвечающее граничным условиям

$$\frac{dU}{dx}(0;p)=V(p), \quad  U(l;p)=0.$$



Отсюда разложением этого изображения в ряд находится решение задачи \eqref{kogan:eq311} в алгебре свертки:

\begin{multline}

\label{kogan:eq32}

u(x;t)=-av(t)\ast \vartheta (t)\ast \biggl[\delta \Bigl(t-\frac{x}{a}  \Bigr)+ \\

+

\sum\limits_{k=1}^\infty (-1)^k\left( 

\delta \Bigl( t-\frac{x}{a}-2\frac{kl}{a} \Bigr)+

\delta \Bigl( t+\frac{x}{a}-2\frac{kl}{a} \Bigr) \right)  

\biggr].

\end{multline}



Выпишем решения, соответствующие таким же частным случаям функции $v(t)$, как и в 

первой задаче: 



\begin{itemize}



\item[а)] если $v(t)$ "--- локально интегрируемая функция, то, используя прежнее обозначение 

$V(z)$, выпишем

$$

u(x;t)=-aV\Bigl(t-\frac{x}{a} \Bigr)+

a\sum\limits_{k=1}^\infty 

(-1)^k\Bigl[

V\Bigl(t-\frac{x}{a}-2\frac{kl}{a}\Bigr)+

V\Bigl(t+\frac{x}{a}-2\frac{kl}{a}\Bigr)\Bigr];

$$



\item[б)] если  $v(t)=\delta (t)$, то

$$

u(x;t)=-a\vartheta \Bigl(t-\frac{x}{a} \Bigr)+

a\sum\limits_{k=1}^\infty (-1)^k

\Bigl[

\vartheta \Bigl( t-\frac{x}{a}-2\frac{kl}{a}\Bigr)+

\vartheta \Bigl( t+\frac{x}{a}-2\frac{kl}{a}\Bigr) \Bigr].

$$



\end{itemize}



Приведем %интересный 

результат непосредственного использования соотношений \eqref{kogan:eq29}. 

Рассмотрим уравнение \eqref{kogan:eq30} при нулевых начальных условиях \eqref{kogan:eq3} и~следующих граничных условиях для закрепленной и свободной на конце $x=l$:

струны~\cite{kogan:bib1}:

\begin{gather}

\label{kogan:eq321}

u(0;t)=g(t), \quad u(l;t)=0;\\

\label{kogan:eq322}

u(0;t)=g(t), \quad u_x (l;t)=0.

\end{gather}

Учитывая \eqref{kogan:eq29}, решения операторного уравнения при граничных условиях

$$

U(0;p)=G(P),\quad U(l;p)=0 \quad  \text{(условия \eqref{kogan:eq321})}

$$ и 

$$U(0;p)=G(P),\quad \frac{dU(l;p)}{dx}=0 \quad  \text{(условия \eqref{kogan:eq322})}

$$ представляем  в виде

$$

U(x;p)=\frac{G(p)(e^{-\frac{x}{a}p}-e^{-(\frac{2l-x}{a}

)p} )}{1-e^{-\frac{2l}{a}p}} \quad \text{и} \quad 

U(x;p)=\frac{G(p) ( e^{-\frac{x}{a}p}+e^{-( \frac{2l-x}{a} )p} 

)}{1+e^{-\frac{2l}{a}p}}.

$$

Отсюда находим решение задач с граничными условиями \eqref{kogan:eq321} и \eqref{kogan:eq322}:

\begin{gather}

\label{kogan:eq33}

u(x;t)=g\Bigl(t-\frac{x}{a}\Bigr)+\sum\limits_{k=1}^\infty 

\Bigl[  

g\Bigl( t-\frac{x}{a}-2\frac{kl}{a}\Bigr)-

g\Bigl( t+\frac{x}{a}-2\frac{kl}{a}\Bigr) \Bigr];\\

\label{kogan:eq34}

u(x;t)=g\Bigl(t-\frac{x}{a} \Bigr)+

\sum\limits_{k=1}^\infty (-1)^k 

\Bigl[

g\Bigl(t-\frac{x}{a}-2\frac{kl}{a} \Bigr)-

g\Bigl(t+\frac{x}{a}-2\frac{kl}{a}\Bigr)\Bigr].

\end{gather}



\smallskip 







{\bf 2.2. Интеграл Дюамеля с точки зрения теории обобщенных функций \cite{kogan:bib1,kogan:bib16,kogan:bib20}.} Введем решение $U(x;t)$, отвечающее однородному уравнению \eqref{kogan:eq1}, \eqref{kogan:eq2}, 

нулевым начальным условиям Коши \eqref{kogan:eq3} и, в зависимости от заданных граничных 

условий \eqref{kogan:eq4} или \eqref{kogan:eq5}, специальным граничным условиям: 

$$

U(0;t)=\vartheta (t), \quad  \lim \limits_{x\to \infty } U(x;t)=0

\quad \text{(граничные условия (\ref{kogan:eq4}))}

$$  или 

$$

U_x  (0;t)=\vartheta (t),\quad \lim \limits_{x\to \infty } U(x;t)=0

\quad \text{(граничные условия \eqref{kogan:eq5})}.$$ 

Для того чтобы полученные далее выражения имели один и тот же вид для 

заданных условий \eqref{kogan:eq4} или \eqref{kogan:eq5}, примем здесь, что в этих условиях обозначения 

правых частей одинаковые, т.~е. $g(t)=v(t)$.



\smallskip



\hypertarget{kogan:th3}{}



\begin{theorem}[3 (интеграл Дюамеля)]Решение однородной задачи {\rm \eqref{kogan:eq1}, \eqref{kogan:eq2}} при нулевых начальных условиях \eqref{kogan:eq3} и граничных условиях \eqref{kogan:eq4} или \eqref{kogan:eq5} можно выразить через соответствующую функцию $U(x;t)$ в алгебре свертки по формуле

\begin{equation}

\label{kogan:eq35}

u(x;t)=\frac{\partial }{\partial t}\bigl(U(x;t)\ast g(t)\bigr)=U_t (x;t)\ast g(t)=U(x;t)\ast g'(t),

\end{equation}

где дифференцирование понимается в смысле теории обобщенных функций.

\end{theorem}



\smallskip



\begin{proof} Будем рассматривать решение в алгебре $D'_+(t)$, при этом переменная $x$ рассматривается 

как параметр. 

В этой алгебре решение представляется в виде 

\begin{equation}

\label{kogan:eq36}

u(x;t)=U(x;t)\ast \delta'(t)\ast g(t).

\end{equation}



Действительно, из определения $U(x;t)$, а также правила дифференцирования 

свертки функций, согласно которому эта операция может быть отнесена только к~$U(x;t)$, 

решение $u(x;t)$ удовлетворяет однородному уравнению \eqref{kogan:eq1}, 

\eqref{kogan:eq2}, нулевым начальным условиям \eqref{kogan:eq3} и граничному условию $u(\infty ;t)=0$. 



Проверим в случае задания граничное условие \eqref{kogan:eq4} для $x=0$: 

$$u(0;t)=U(0;t)\ast \delta'(t)\ast g(t)=\vartheta (t)\ast \delta'(t)\ast 

g(t)=g(t).$$ 



Проверка граничного условия \eqref{kogan:eq5} проводится аналогично. 



Из \eqref{kogan:eq36} следует, что 

$$u(x;t)=\big( U(x;t)\ast g(t) \big)\ast \delta'

(t)=\frac{\partial }{\partial t} \bigl( U(x;t)\ast g(t)\bigr),$$ 

откуда 

следует~\eqref{kogan:eq35}.

\end{proof}



\smallskip



Приведем некоторые практические формулы реализации теоремы Дюамеля, которые 

нетрудно получить из~\eqref{kogan:eq35} в зависимости от заданной функции $g(t)$ в~начальных условиях~\cite{kogan:bib20}:

\begin{multline}

\label{kogan:eq37}

u(x;t)=\frac{\partial }{\partial t}\int_0^t U(x;t-\tau )g(\tau )d\tau 

=\int_0^t U_t (x;t-\tau )g(\tau )d\tau =\\

=U(x;t)g(0)+\int _0^t U(x;t-\tau )g'(\tau )d\tau,

\end{multline}

где $g(t)$ и $g'(t)$ "--- обычные кусочно"=непрерывные функции;

\begin{equation}

\label{kogan:eq38}

u(x;t)=\sum\limits_{i=1}^m c_i \vartheta \big(t-t_i\big )U\big(x;t-t_i\big),

\end{equation}

где $g(t)=\sum\limits_{i=1}^m c_i \vartheta (t-t_i)$, $c_i$ "--- 

постоянные.



\smallskip



\begin{remark}[2] Теорема \hyperlink{kogan:th3}{3}, как нетрудно видеть, справедлива, если однородное 

уравнение \eqref{kogan:eq1}, \eqref{kogan:eq2}, отвечающее нулевым начальным условиям \eqref{kogan:eq3}, задано на 

отрезке $x\in [0;l]$. 

В этом случае должны быть указаны  граничные условия при $x=l$, которым в месте с условием $U(0;t)=\vartheta 

(t)$ или $U_x (0;t)=\vartheta (t)$ должна удовлетворять функция $U(x;t)$.

\end{remark}



\smallskip



На известном примере покажем применение теоремы.



\smallskip



{\bf 2.3. Распределение температуры в полуограниченном стержне.} \linebreak

Пусть в начальный момент времени тонкий стержень имеет нулевую 

температуру и задана функция изменения температуры "--- $g(t)$ при 

$x=0$. 

Таким образом, решаем  уравнение теплопроводности при  следующих начальных и~граничных условиях: 

\begin{equation}

\label{kogan:eq39}

\Bigl(\frac{\partial }{\partial t}-a^2\frac{\partial ^2}{\partial x^2} 

\Bigr)u(x;t)=0;

\quad

u(x;0)=0;

\quad

u(0;t)=g(t);

\quad

\lim \limits_{x\to \infty } u(x;t)=0.

\end{equation}



Здесь удобно применить теорему Дюамеля. Найдем вспомогательное решение 

$U(x;t)$, отвечающее начальному условию $U(0;t)=\vartheta (t)$. 

Соответствующее операторное уравнение 

$$pU-a^2\frac{d^2U}{dx^2}=0$$

при граничных условиях 

$$U (0;p)= \frac 1 p,  \quad \lim \limits_{x\to \infty, \mathop{\rm Re} p>0} U(x;p)\to 

0$$ имеет решение 

$$U(x;p)=\frac{e^{-\frac{x}{a}\sqrt p }}{p}.$$ 

Используя известное соотношение 

преобразования Лапласа~\cite{kogan:bib3} 

$$\frac{e^{-\alpha \sqrt p}}{p}\leftrightarrow 

{\rm Erf}\Big( \frac{\alpha} { 2\sqrt t} \Bigr)\quad  (\alpha >0),$$ получаем 

\begin{equation}

\label{kojgan:eq40}

U(x;t)={\rm Erf}\Bigl(\frac x {2a\sqrt t} \Bigr)=1-\mathop{\rm erf} \Bigl(\frac x {2a\sqrt t} \Bigr),

\end{equation}

где $$\mathop{\rm erf} ( z )=\frac{2}{\sqrt \pi }\int _0^z e^{-\tau ^2}d\tau.$$



Выпишем решение задачи  \eqref{kogan:eq39} для частных случаев:

\begin{itemize}

 

\item [а)] если $g(t)$ "--- непрерывная функция, то, используя второе соотношение формулы 

\eqref{kogan:eq35} или \eqref{kogan:eq37}, будем иметь

\[

u(x;t)=\int_0^t U_t (x;t-\tau )g(\tau )d\tau  =\frac{x}{2a\sqrt \pi 

}\int _0^t \frac{g(\tau )}{(t-\tau )^{3/2 }} e^{-\frac{x^2}{4a^2(t-\tau )}}d\tau;

\]



\item [б)] если $g(t)=\vartheta (t-1)-\vartheta (t-2)$, то, используя третье 

соотношение формулы \eqref{kogan:eq35} или формулу \eqref{kogan:eq38}, будем иметь

\[

u(x;t)=\vartheta (t-1)\mathop{\rm Erf}\Bigl( \frac{x} {2a \sqrt {t-1} }\Bigr)-\vartheta (t-2)

\mathop{\rm Erf} \Bigl(\frac x {2a\sqrt {t-2}} \Bigr).

\]





\end{itemize}



\smallskip



\Section[n]{Выводы} 

Приведенная теория и практика применения операторного 

исчисления Микусинского на основе алгебры свертки обобщенных функций $D'_+$ и $D'_-$ к решению линейных уравнений в~частных производных с~постоянными коэффициентами показывает следующее:



\begin{itemize}

\item[1)] возможно рассмотрение уравнений, заданных на всей оси; 



\item[2)] использование формальных свойств алгебры свертки и правил 

дифференцирования обобщенных функций позволяет получить компактные 

доказательства теорем, а также упростить нахождение и форму представления 

решений, в том числе при задании разрывных начальных условий;



\item[3)] применение рядов сверточных операторов обобщенных функций пополняет 

арсенал средств для нахождения решений "--- оригиналов по найденному выражению 

изображения. 





\end{itemize}



К ограничению применения рассматриваемого операторного исчисления следует 

отнести то, что он использует функции только одной переменной. 

В~то же время 

вместе с более общим преобразованием Фурье и другими интегральными 

преобразованиями благодаря большой базе данных для преобразования Лапласа 

предлагаемый подход может найти место в математической физике. Это, прежде 

всего, традиционное для операционного исчисления направление "--- {\sl решение 

нестационарных задач}.



\smallskip



\AdditionalContent{Конкурирующие интересы}{Конкурирующих интересов не имею.}



\AdditionalContent{Авторский вклад и ответственность}{Я несу полную ответственность за предоставление окончательной версии рукописи в печать. Окончательная версия 

рукописи мною одобрена.}



\AdditionalContent{Финансирование}{Исследование выполнялось без финансирования.}



%\AdditionalContent{Благодарность}{Авторы благодарны рецензенту за тщательное прочтение статьи и~ценные предложения и комментарии.}



\vspace{-3mm}



\begin{thebibliography}{99}



\Bibitem{kogan:bib1}

\by Courant~R., Hilbert~D.

\book Methods of Mathematical Physics

\bookvol 2

\volinfo Partial Differential Equations

\publaddr New York

\publ Interscience Publ.

\yr 1962

\totalpages xxii+830



\Bibitem{kogan:bib2}

\by Doetsch~G.

\book Introduction to the Theory and Application of the Laplace Transformation

\publ Springer-Verlag 

\publaddr Berlin, Heidelberg

\yr 1974

\totalpages viii+327~pp \nofrills

\crossref{10.1007/978-3-642-65690-3}



\RBibitem{kogan:bib3}

\by Лаврентьев~М.~А., Шабат~Б.~В.

\book Методы теории функций комплексного переменного

\publaddr М.

\publ Наука

\yr 1987

\totalpages 688

\zmath{https://zbmath.org/?q=an:04029882}

%\misctransl

%\lang In Russian

%\by Lavrent'ev~M.~A., Shabat~B.~V.

%\book Methods of the theory of functions in a complex variable

%\publaddr Moscow

%\publ Nauka

%\yr 1987

%\totalpages 688



\RBibitem{kogan:bib4}

\by Федорюк~М.~В.

\paper Интегральные преобразования

\inbook Анализ~--~1

\serial Итоги науки и техн. Сер. Соврем. пробл. мат. Фундам. направления

\yr 1986

\vol 13

\pages 211--253

\publ ВИНИТИ

\publaddr М.

\mathnet{http://mi.mathnet.ru/intf74}

\mathscinet{http://www.ams.org/mathscinet-getitem?mr=899754}

\zmath{https://zbmath.org/?q=an:0658.44001}

%\misctransl

%\lang In Russian

%\by Fedoryuk~M.~V.

%\paper Integral Transforms

%\inbook Analysis

%\vol  1

%\serial Integral Representations and Asymptotic Methods

%\ed Gamkrelidze~R~.V.

%\publ Springer -- Verlag

%\publaddr Berlin Heidelberg

%\yr 1989

%\pages 193--232



\Bibitem{kogan:bib5}

\by Sharma~J.~N., Singh~K.

\book Partial Differential Equations for Engineers and Scientists

\publaddr New Delhi

\publ Narosa Publishing House

\yr 2011

\totalpages 354 

\zmath{https://zbmath.org/?q=an:05842521}



\RBibitem{kogan:bib6}

\by Волков~И.~К., Канатников~А.~Н.

\book Интегральные преобразования и операционное исчисление

\publaddr М.

\eds В.~С.~Зарубин, А.~П.~Крищенко 

\publ МГТУ им.~Н.~Э.~Баумана

\yr 2015

\totalpages 227

\serial Математика в техническом университете 

\vol 11

%\misctransl

%\lang In Russian

%\by Volkov~I.~K., Kanatnikov~A.N.

%\book Integral Transforms and Operational Calculus

%\publaddr Moscow

%\publ Bauman Moscow State Technical University

%\yr 2015

%\totalpages 227



\Bibitem{kogan:bib7}

\zmath{https://zbmath.org/?q=an:03074402}

\by Schwartz~L.

\paper Transformation de Laplace des distributions

\jour Meddel. Lunds Univ. Mat. Sem. Suppl.-band M. Riesz

\yr 1952

\pages 196--206 

\lang In French





\Bibitem{kogan:bib8}

\zmath{https://zbmath.org/?q=an:03081985}

\lang In French

\by Lions~J.~L.

\paper Supports dans la transformation de Laplace

\jour J. Anal. Math. 

\vol 2

\pages 369--380 

\yr 1953



\Bibitem{kogan:bib9}

\zmath{https://zbmath.org/?q=an:03166049}

\by Schwartz~L.

\book M\'ethodes math\'ematiques pour les sciences physiques. Avec le concours de Denise Huet

\serial Enseign. des sciences

\publaddr Paris

\publ Hermann \& Cie

\totalpages 392 

\yr 1961

\lang In French





\RBibitem{kogan:bib10}

\by Владимиров~В.~С.

\book Обобщенные функции в математической физике

\publaddr М.

\publ Наука

\yr 1979

\totalpages 319

\zmath{https://zbmath.org/?q=an:0515.46033}

%\misctransl

%\lang In Russian

%\by Vladimirov~V.~S.

%\book Generalized Functions in Mathematical Physics

%\publaddr  Moscow

%\publ Mir Publishers

%\yr 1979

%\totalpages 362



\RBibitem{kogan:bib11}

\by Владимиров~В.~С.

\book Уравнения математической физики

\publaddr М.

\publ Наука

\yr 1981

\totalpages 512

%\misctransl

%\lang In Russian

%\by Vladimirov~V.~S.

%\book Equations of Mathematical Physics. 4 th edition

%\publaddr Moscow

%\publ Mir Publishers

%\yr 1984

%\totalpages 464



\RBibitem{kogan:bib12}

\by Брычков~Ю.~А., Прудников~А.~П.

\book Интегральные преобразования обобщенных функций

\publaddr М.

\publ Наука

\yr 1977

\totalpages 287

\zmath{https://zbmath.org/?q=an:0464.46039}

%\misctransl

%\lang In Russian

%\by Brychkov~Y.~A., Prudnikov~A.~P.

%\book Integral transformations of generalized functions

%\publaddr Moscow

%\publ Nauka

%\yr 1977

%\totalpages 288



\Bibitem{kogan:bib13}

\by Kecs~W., Teodorescu~P.~P.

\book Application of the distribution theory in the mechanics

\lang In Romanian

\publaddr Bucuresti

\publ Editura Academiei Republicii Socialiste Romania

\totalpages 438 

\yr 1970

\zmath{https://zbmath.org/?q=an:03352831}



\Bibitem{kogan:bib14}

\by Mikusi\'nski~J.

\book Operational calculus

\serial Internat. Series of Monographs on Pure and Applied Mathematics

\vol 8

\publaddr New York

\publ Pergamon Press

\yr 1959

\totalpages 495

\zmath{https://zbmath.org/?q=an:0088.33002}



\RBibitem{kogan:bib15}

\by Коган~И.~Л.

\paper Построение операторного исчисления Микусинского на основе 

алгебры свертки обобщенных функций. Основные положения 

\jour  Вестн. Сам. гос. техн. ун-та. Сер. Физ.-мат. науки

\yr 2012

\issue 2(27)

\pages 44--52

\mathnet{http://mi.mathnet.ru/vsgtu1013}

\crossref{10.14498/vsgtu1013}



\RBibitem{kogan:bib16}

\by Коган~И.~Л.

\paper Построение операторного исчисления Микусинского на основе 

алгебры свертки обобщенных функций. Теоремы и начало применения

\jour  Вестн. Сам. гос. техн. ун-та. Сер. Физ.-мат. науки

\yr 2013

\issue 3(32) 

\pages 56--68

\mathnet{http://mi.mathnet.ru/vsgtu1119}

\crossref{10.14498/vsgtu1119}





\Bibitem{kogan:bib17}

\by H\"{o}rmander~L.

\book Linear partial differential operators

\serial Die Grundlehren der mathematischen Wissenschaften

\vol 116

\publaddr Berlin,  G\"{o}ttingen, Heidelberg

\publ Springer-Verlag

\totalpages vii+285 

\yr 1963

\zmath{https://zbmath.org/?q=an:0108.09301}



\RBibitem{kogan:bib18}

\by Гельфанд~И.~М., Шилов~Г.~Е.

\book Обобщенные функции

\bookvol 1

\volinfo Обобщенные функции и действия над ними

\publaddr М.

\publ Физматлит

\yr 1959

\totalpages 470

\zmath{https://zbmath.org/?q=an:0091.11102}

%\misctransl

%\lang In Russian

%\by Gelfand~I.~M., Shilov~G.~E.

%\book Generalized functions

%\vol 1. Properties and operations

%\publaddr New York

%\publ Academic Press

%\yr 1964

%\totalpages 423



\Bibitem{kogan:bib19}

\by Korn~G.~A., Korn~T.~M.

\book Mathematical handbook for scientists and engineers: definitions, theorems, and formulas for reference and review

\bookinfo Reprint with corrections of the 2nd revised and enlarged edition 1968. 

\serial Dover Civil and Mechanical Engineering

\publaddr Mineola, NY

\publ Dover Publications 

\totalpages xvii+1130 

\yr 2003

\zmath{https://zbmath.org/?q=an:1326.00020}



\RBibitem{kogan:bib20}

\by Коган~И.~Л.

\paper Метод интеграла Дюамеля для обыкновенных дифференциальных уравнений с постоянными коэффициентами с~точки зрения теории обобщенных функций

\jour Вестн. Сам. гос. техн. ун-та. Сер. Физ.-мат. науки

\yr 2010

\issue 1(20)

\pages 37--45

\mathnet{http://mi.mathnet.ru/vsgtu673}

\crossref{10.14498/vsgtu673}





\end{thebibliography}



\renewcommand{\baselinestretch}{0.85}



\newpage



\vspace{5mm}



\makeentitle



\renewcommand{\baselinestretch}{0.93}







\AdditionalContent{Competing interests}{I declare that I have no competing interests} 

\AdditionalContent{Author's Responsibilities}{I take full responsibility for submitting the final manuscript in print. I approved the final version of the manuscript.} 

\AdditionalContent{Funding}{The research has not had any funding.}









\begin{thebibliography}{99}



\Bibitem{kogan:bib1:en}

\by Courant~R., Hilbert~D.

\book Methods of Mathematical Physics

\bookvol 2

\volinfo Partial Differential Equations

\publaddr New York

\publ Interscience Publ.

\yr 1962

\totalpages xxii+830



\Bibitem{kogan:bib2:en}

\by Doetsch~G.

\book Introduction to the Theory and Application of the Laplace Transformation

\publ Springer-Verlag 

\publaddr Berlin, Heidelberg

\yr 1974

\totalpages viii+327~pp \nofrills

\crossref{10.1007/978-3-642-65690-3}



\Bibitem{kogan:bib3:en}

\lang In Russian 

\by Lavrent'ev~M.~A., Shabat~B.~V.

\book Metody teorii funktsii kompleksnogo peremennogo \rm [Methods for the theory of functions of a complex variable]

\publaddr Moscow

\publ Nauka

\yr 1987

\totalpages 688





\Bibitem{kogan:bib4:en}

\by Fedoryuk~M.~V.

\paper Integral Transforms

\inbook Analysis~I. Integral Representations and Asymptotic Methods

\serial Encyclopaedia of Mathematical Sciences

\vol 13

\ed R.~V.~Gamkrelidze

\publ Springer-Verlag

\publaddr Berlin, Heidelberg

\yr 1989

\pages 193--232

\crossref{10.1007/978-3-642-61310-4_2}



\Bibitem{kogan:bib5:en}

\by Sharma~J.~N., Singh~K.

\book Partial Differential Equations for Engineers and Scientists

\publaddr New Delhi

\publ Narosa Publishing House

\yr 2011

\totalpages 354 

\zmath{https://zbmath.org/?q=an:05842521}



\Bibitem{kogan:bib6:en}

\lang In Russian

\by Volkov~I.~K., Kanatnikov~A.~N.

\book Integral'nye preobrazovaniia i operatsionnoe ischislenie \rm [Integral transforms and operational calculus]

\publaddr Moscow

\eds V.~S.~Zarubin, A.~P.~Krishchenko 

\publ Bauman Moscow State Techn. Univ.

\yr 2015

\totalpages 227

\serial Matematika v tekhnicheskom universitete \rm [Mathe\-matics in the Technical University]

\vol 11





\Bibitem{kogan:bib7:en}

\zmath{https://zbmath.org/?q=an:03074402}

\by Schwartz~L.

\paper Transformation de Laplace des distributions

\jour Meddel. Lunds Univ. Mat. Sem. Suppl.-band M. Riesz

\yr 1952

\pages 196--206 

\lang In French





\Bibitem{kogan:bib8:en}

\zmath{https://zbmath.org/?q=an:03081985}

\lang In French

\by Lions~J.~L.

\paper Supports dans la transformation de Laplace

\jour J. Anal. Math. 

\vol 2

\pages 369--380 

\yr 1953



\Bibitem{kogan:bib9:en}

\zmath{https://zbmath.org/?q=an:03166049}

\by Schwartz~L.

\book M\'ethodes math\'ematiques pour les sciences physiques. Avec le concours de Denise Huet

\serial Enseign. des sciences

\publaddr Paris

\publ Hermann \& Cie

\totalpages 392 

\yr 1961

\lang In French





\Bibitem{kogan:bib10:en}

\lang In Russian

\by Vladimirov~V.~S.

\book Obobshchennye funktsii v matematicheskoi fizike \rm [Generalized Functions in Mathematical Physics]

\publaddr Moscow

\publ Nauka

\yr 1979

\totalpages 319





\Bibitem{kogan:bib11:en}

\lang In Russian

\by Vladimirov~V.~S.

\book Uravneniia matematicheskoi fiziki \rm [Equations of Mathematical Physics]

\publaddr Moscow

\publ Nauka

\yr 1981

\totalpages 512





\Bibitem{kogan:bib12:en}

\lang In Russian

\by Brychkov~Yu.~A., Prudnikov~A.~P.

\book Integral'nye preobrazovaniia obobshchennykh funktsii \rm [Integral transformations of generalized functions]

\publaddr Moscow

\publ Nauka

\yr 1977

\totalpages 287



\Bibitem{kogan:bib13:en}

\by Kecs~W., Teodorescu~P.~P.

\book Application of the distribution theory in the mechanics

\lang In Romanian

\publaddr Bucuresti

\publ Editura Academiei Republicii Socialiste Romania

\totalpages 438 

\yr 1970

\zmath{https://zbmath.org/?q=an:03352831}



\Bibitem{kogan:bib14:en}

\by Mikusi\'nski~J.

\book Operational calculus

\serial Internat. Series of Monographs on Pure and Applied Mathematics

\vol 8

\publaddr New York

\publ Pergamon Press

\yr 1959

\totalpages 495

\zmath{https://zbmath.org/?q=an:0088.33002}



\Bibitem{kogan:bib15:en}

\lang In Russian

\by Kogan~I.~L.

\paper Construction of Mikusinski operational calculus based on the convolution algebra of distributions. Basic provisions

\jour Vestn. Samar. Gos. Tekhn. Univ., Ser. Fiz.-Mat. Nauki \rm [J. Samara State Tech. Univ., Ser. Phys. Math. Sci.]

\yr 2012

\issue 2(27)

\pages 44--52

\mathnet{http://mi.mathnet.ru/vsgtu1013}

\crossref{10.14498/vsgtu1013}



\Bibitem{kogan:bib16:en}

\by Kogan~I.~L.

\paper Construction of Mikusinski operational calculus based on the convolution algebra of distributions. The theorems and the beginning of use

\jour Vestn. Samar. Gos. Tekhn. Univ., Ser. Fiz.-Mat. Nauki [J. Samara State Tech. Univ., Ser. Phys. Math. Sci.]

\yr 2013

\issue 3(32) 

\pages 56--68

\mathnet{http://mi.mathnet.ru/vsgtu1119}

\crossref{10.14498/vsgtu1119}





\Bibitem{kogan:bib17:en}

\by H\"{o}rmander~L.

\book Linear partial differential operators

\serial Die Grundlehren der mathematischen Wissenschaften

\vol 116

\publaddr Berlin,  G\"{o}ttingen, Heidelberg

\publ Springer-Verlag

\totalpages vii+285 

\yr 1963

\zmath{https://zbmath.org/?q=an:0108.09301}



\Bibitem{kogan:bib18:en}

\by Gelfand~I.~M., Shilov~G.~E.

\book Generalized functions

\bookvol 1

\volinfo  Properties and operations

\publaddr New York, London

\publ Academic Press

\yr 1964

\totalpages xviii+423

\zmath{https://zbmath.org/?q=an:0115.33101}





\Bibitem{kogan:bib19:en}

\by Korn~G.~A., Korn~T.~M.

\book Mathematical handbook for scientists and engineers: definitions, theorems, and formulas for reference and review

\bookinfo Reprint with corrections of the 2nd revised and enlarged edition 1968. 

\serial Dover Civil and Mechanical Engineering

\publaddr Mineola, NY

\publ Dover Publications 

\totalpages xvii+1130 

\yr 2003

\zmath{https://zbmath.org/?q=an:1326.00020}



\Bibitem{kogan:bib20:en}

\lang In Russian

\by Kogan~I.~L.

\paper Method of Duhamel Integral for Ordinary Differential Equations with Constant Coefficients in Respect to the Theory of Distributions

\jour Vestn. Samar. Gos. Tekhn. Univ., Ser. Fiz.-Mat. Nauki \rm [J. Samara State Tech. Univ., Ser. Phys. Math. Sci.]

\yr 2010

\issue 1(20)

\pages 37--45

\mathnet{http://mi.mathnet.ru/vsgtu673}

\crossref{10.14498/vsgtu673}





\end{thebibliography}





\renewcommand{\baselinestretch}{0.9}



\newpage



%\end{document}

